\documentclass{article}
\usepackage[margin=1in]{geometry}
\usepackage{amsmath}
\usepackage{amssymb}
\usepackage{enumitem}
\setlength{\parindent}{0pt}

\begin{document}

\begin{center}
{\LARGE\bfseries Factors, Multiples and Primes: Answer Key}
\end{center}

\section*{Factors}
\textbf{True/false:} 1 T, 2 T, 3 F, 4 F, 5 F.

\begin{enumerate}[label=\arabic*., nosep, leftmargin=*]
\item \(12\): \(1,2,3,4,6,12\).
\item \(20\): \(1,2,4,5,10,20\).
\item \(36\): \(1,2,3,4,6,9,12,18,36\).
\item \(16\) has five factors: \(1,2,4,8,16\).
\item \(\mathrm{HCF}(12,18)=6\).
\item \(\mathrm{HCF}(24,36)=12\).
\item \(1\times48,\;2\times24,\;3\times16,\;4\times12,\;6\times8\).
\item \(24\), with eight factors: \(1,2,3,4,6,8,12,24\).
\item Any square of a prime, e.g.\ \(4,9,25,\dots\)
\item Factors come in pairs, but in a square number the square root pairs with itself, leaving one factor unpaired; hence the number of factors is odd.
\end{enumerate}

\section*{Multiples}
\textbf{True/false:} 1 T, 2 T, 3 F, 4 F, 5 F.

\begin{enumerate}[label=\arabic*., nosep, leftmargin=*]
\item \(4,8,12,16,20\).
\item \(7,14,21,28,35\).
\item \(24,30,36,42,48\).
\item Yes; the digit sum is \(9\), and \(126=9\times14\).
\item \(\mathrm{LCM}(4,6)=12\).
\item \(\mathrm{LCM}(8,12)=24\).
\item \(\mathrm{LCM}(5,6,9)=90\).
\item \(20\).
\item In \(36\) seconds.
\item \(61\).
\end{enumerate}

\section*{Primes}
\textbf{True/false:} 1 F, 2 T, 3 F, 4 T, 5 F.

\begin{enumerate}[label=\arabic*., nosep, leftmargin=*]
\item \(2,3,5,7\).
\item \(11,13,17,19\).
\item No; \(27=3\times9\).
\item \(31\).
\item No; \(91=7\times13\).
\item Four: \(23,29,31,37\).
\item No; \(143=11\times13\).
\item \(5+19,\;7+17,\;11+13\).
\item Test \(2,3,5,7,11,13\); none divides \(211\), so \(211\) is prime.
\item \(n^2-1=(n-1)(n+1)\), and both factors are greater than \(1\) when \(n>2\).
\end{enumerate}

\section*{Is it prime?}
The primes from \(1\) to \(100\) are
\[
2,3,5,7,11,13,17,19,23,29,31,37,41,43,47,53,59,61,67,71,73,79,83,89,97.
\]
All other whole numbers from \(1\) to \(100\) are not prime.

\section*{Prime factor decomposition}
\textbf{True/false:} 1 T, 2 F, 3 T, 4 F, 5 F.

\begin{enumerate}[label=\arabic*., nosep, leftmargin=*]
\item \(12=2^2\times3\).
\item \(18=2\times3^2\).
\item \(60=2^2\times3\times5\).
\item \(84=2^2\times3\times7\).
\item \(200=2^3\times5^2\).
\item \(360=2^3\times3^2\times5\).
\item \(1001=7\times11\times13\).
\item \(600\).
\item \(2n=2^4\times3\times5\).
\item \(5\times3\times2=30\) factors.
\end{enumerate}

\section*{HCF and LCM from lists}
\begin{enumerate}[label=\arabic*., nosep, leftmargin=*]
\item \(4\) and \(6\): \(\mathrm{HCF}=2,\ \mathrm{LCM}=12\).
\item \(6\) and \(8\): \(\mathrm{HCF}=2,\ \mathrm{LCM}=24\).
\item \(6\) and \(9\): \(\mathrm{HCF}=3,\ \mathrm{LCM}=18\).
\item \(8\) and \(12\): \(\mathrm{HCF}=4,\ \mathrm{LCM}=24\).
\item \(10\) and \(15\): \(\mathrm{HCF}=5,\ \mathrm{LCM}=30\).
\item \(9\) and \(12\): \(\mathrm{HCF}=3,\ \mathrm{LCM}=36\).
\item \(12\) and \(18\): \(\mathrm{HCF}=6,\ \mathrm{LCM}=36\).
\item \(14\) and \(21\): \(\mathrm{HCF}=7,\ \mathrm{LCM}=42\).
\item \(15\) and \(20\): \(\mathrm{HCF}=5,\ \mathrm{LCM}=60\).
\item \(8\) and \(9\): \(\mathrm{HCF}=1,\ \mathrm{LCM}=72\).
\end{enumerate}

\section*{HCF and LCM by prime factor decomposition}
\textbf{True/false:} 1 F, 2 T, 3 T, 4 T, 5 F.

\begin{enumerate}[label=\arabic*., nosep, leftmargin=*]
\item \(\mathrm{HCF}(12,18)=6\).
\item \(\mathrm{LCM}(12,18)=36\).
\item \(\mathrm{HCF}(24,60)=12\).
\item \(\mathrm{LCM}(24,60)=120\).
\item \(\mathrm{HCF}(84,120)=12,\ \mathrm{LCM}(84,120)=840\).
\item \(\mathrm{HCF}(126,180)=18,\ \mathrm{LCM}(126,180)=1260\).
\item \(\mathrm{HCF}=2^2\times3\times5,\ \mathrm{LCM}=2^3\times3^2\times5^2\).
\item The other number is \(18\).
\item \(\mathrm{HCF}=2^3\times3^2,\ \mathrm{LCM}=2^5\times3^4\times7\times11\).
\item \(12\) and \(180\), or \(36\) and \(60\); both pairs work.
\end{enumerate}

\section*{Square numbers from prime factors}
\textbf{True/false:} 1 T, 2 F, 3 T, 4 F, 5 T.

\begin{enumerate}[label=\arabic*., nosep, leftmargin=*]
\item Yes; \(36=2^2\times3^2=(2\times3)^2=6^2\).
\item Yes; the square root is \(2\times3\times5=30\).
\item No; the power of \(2\) is odd.
\item \(72\).
\item Yes; \(400=2^4\times5^2=(2^2\times5)^2=20^2\).
\item No; \(150=2\times3\times5^2\), and the powers of \(2\) and \(3\) are odd.
\item \(1200\).
\item \(k=2\times3=6\).
\item Multiply by \(2\times7=14\).
\item Divide by \(2\times3\times5=30\), giving \(36=6^2\).
\end{enumerate}

\section*{Problem solving with HCF and LCM}
\textbf{True/false:} 1 T, 2 F, 3 T, 4 T, 5 F.

\begin{enumerate}[label=\arabic*., nosep, leftmargin=*]
\item In \(12\) minutes.
\item Each piece is \(6\) cm.
\item \(12\) packs, each with \(3\) pens and \(4\) pencils.
\item At \(8{:}01{:}30\) pm.
\item A \(36\) cm by \(36\) cm square.
\item At \(10{:}12\).
\item Small cog: \(4\) turns; large cog: \(3\) turns.
\item Tiles are \(72\) cm square; \(35\) tiles are needed.
\item After \(6\) minutes.
\item \(120\) hot dogs: \(12\) packs of hot dogs, \(15\) packs of buns, \(10\) packs of sauce sachets.
\end{enumerate}

\end{document}